\unnumberedsection{Further Readings}

\begin{referencelist}
  \refitem Burgis, T. (2020). \emph{Kleptopia: How dirty money is conquering the world}. William Collins. Follows specific flows of illicit money across borders and shows how secrecy structures connect crime, business, and politics.
  \refitem Financial Action Task Force. (2023). \emph{International standards on combating money laundering and the financing of terrorism and proliferation: The FATF recommendations}. FATF. The global rulebook on anti-money-laundering and beneficial ownership; essential primary reading for drafting workable mechanisms. \url{https://www.fatf-gafi.org/en/publications/Fatfrecommendations/Fatf-recommendations.html}
  \refitem INTERPOL. (2026). \emph{Global financial fraud threat assessment} (2nd ed.). The most current institutional picture of transnational financial fraud, useful for grounding debate in real figures and methods. \url{https://www.interpol.int/Media/Documents/Publications/Financial-Crime/INTERPOL-Global-Financial-Fraud-Threat-Assessment-March-2026}
  \refitem Martha, R. S. J. (2010). \emph{The legal foundations of INTERPOL} (2nd ed.). Hart Publishing. The standard legal treatment of the Organization’s powers, limits, and notice system, for delegates who want the committee’s mechanics in detail.
\end{referencelist}
